\chapter*{Abstract}
\addcontentsline{toc}{chapter}{Abstract}

We document \textsf{KELP}, a Lagrangian particle-tracking engine that
integrates the trajectories of passive particles through the velocity fields of
unstructured-mesh ocean circulation models, with primary support for the FESOM
family --- FESOM2, on purely triangular meshes, and FESOM-C, on mixed
triangular--quadrilateral meshes --- and with readers for ocean fields remapped to the HEALPix grid and
for structured NEMO grids.
This manual is a numerical and implementation reference: for each component of
the engine we state the governing equations, give the discrete formulae exactly
as implemented, illustrate the underlying geometry with schematics, and cite the
location of the corresponding Python source. We treat the engine as the object of
study rather than any particular application; the FESOM2 \emph{core2} global mesh recurs only as a
concrete example.

The material is organized along the path a single time step travels through the
solver: locating a particle on the unstructured mesh, with a proof that the nearest node found is the
nearest; interpolating the velocity in space (horizontally by barycentric or least-squares schemes,
vertically between depth interfaces or along sigma surfaces) and in time (by linear or monotone
piecewise-cubic reconstruction); advancing the position with an explicit or adaptive Runge--Kutta
scheme; and applying boundary conditions. Stochastic diffusion, the particle lifecycle, the
forcing reader, trajectory output, the configuration surface, verification, and
computational cost --- measured across thread layouts, mesh sizes, particle counts and node counts ---
are each treated in turn, and a final part describes the agent skills through which experiments are set
up, run and analysed. We point to the engine's
physics-validation tests --- the analytic-flow benchmarks, the vertical streamfunction
(continuity) conservation tests, the reversibility round-trip, and the checks against real model
output --- by name, as evidence that the
numerics behave as specified, without reproducing their internals.
